\chapter*{Kurzfassung}
\addcontentsline{toc}{chapter}{Kurzfassung}

Generative KI-Systeme werden von Studierenden inzwischen routinemäßig auch bei komplexen Aufgaben
eingesetzt, was die Gefahr birgt, dass Lösungen ungeprüft übernommen werden. Diese Arbeit untersucht,
ob ein speziell für Reflexion gestaltetes KI-System (Reflexionsbot) dazu führt, dass Studierende
Prozessentscheidungen im Software-Prozessmanagement eigenständiger und reflektierter bearbeiten als
mit einem generischen KI-System.

Dazu wurde ein Reflexionsbot umgesetzt, der mittels sokratischem Fragen und adaptiver Antworttiefe
statt direkter Lösungen begleitet. In einem Mini-Experiment in der Lehrveranstaltung \glqq Grundlagen
des Software Engineering 1\grqq{} im Bachelorstudiengang Software Engineering der Hochschule
Heilbronn bearbeitete eine Gruppe zwei realitätsnahe Prozessszenarien mit dem Reflexionsbot, eine
Vergleichsgruppe mit einem generischen KI-System. Die 31 Bearbeitungen und die zugehörigen
Chatverläufe wurden mittels strukturierender qualitativer Inhaltsanalyse nach Hatton und Smith
kodiert, ergänzt um die Autorenschaft des Textes als eigene Kategorie, und mit einem Prä-Post-Survey
(n=22 bzw. n=20) trianguliert.

Der Anteil der Bearbeitungen ohne eindeutig studentische Autorenschaft ist in der
Reflexionsbot-Gruppe geringer (8 von 16 gegenüber 12 von 15). Dieser Abstand ist jedoch teilweise
durch die Aufgabenvorgabe für die Vergleichsgruppe bedingt und fällt bei den gesichert KI-generierten
Fällen deutlich kleiner aus (7 von 16 gegenüber 8 von 15). Die direkten Lernmaße stützen keinen
Vorteil des Reflexionsbots. Im Wissenstest verhindert ein Deckeneffekt jede Aussage, und die gezeigte
Reflexionstiefe liegt deskriptiv sogar niedriger als in der Vergleichsgruppe (mittlere Stufe 0.75
gegenüber 1.33), bei drei gegenüber acht auswertbaren Fällen allerdings ohne Aussagekraft für die
Richtung. Günstiger fällt lediglich die Selbsteinschätzung zum Begründen und Abwägen aus. Didaktisches
Design beeinflusst damit das Nutzungsverhalten, ein Lernvorteil ergibt sich daraus aber nicht
automatisch. Entscheidend ist die Bereitschaft der Studierenden, sich auf den Dialog einzulassen.

\vspace{1em}
\noindent\textbf{Schlagwörter:} Generative KI, Reflexionsbot, Sokratisches Fragen, Software
Engineering Education, Mixed-Methods, Qualitative Inhaltsanalyse
